\par\noindent\textbf{(1)} Pauli向量 \(\boldsymbol h=(\Delta,0,vt)\) 的模为 \(\sqrt{\Delta^2+v^2t^2}\)，故本征能 \(E_\pm=\pm\frac12\sqrt{\Delta^2+v^2t^2}\)。能隙在 \(t=0\) 取最小值 \(\Delta\)，而非零；当 \(\Delta=0\) 才是真交叉。
\par\noindent\textbf{(2)} 令 \(\tan\theta=\Delta/(vt)\)，选择连续支，\(\dot\theta=-v\Delta/(\Delta^2+v^2t^2)\)。绝热基的非对角耦合模为 \(|\dot\theta|/2\)，在交叉中心达到 \(v/(2\Delta)\) 的最大值，解释小能隙和高速扫描容易产生跃迁。
\par\noindent\textbf{(3)} 通过将二分量Schrödinger方程化为抛物柱函数方程，并连接 \(t\to\pm\infty\) 渐近解，可得从初始绝热基态跳至末态激发支的概率 \(P_{\rm LZ}=\exp[-\pi\Delta^2/(2\hbar v)]\)。当 \(v\to0\) 指数抑制跃迁，快扫描 \(v\to\infty\) 则 \(P_{\rm LZ}\to1\)。
\par\noindent\textbf{(4)} 无量纲绝热参数为 \(\delta=\Delta^2/(4\hbar v)\)，指数可写 \(-2\pi\delta\)。该精确公式要求线性时间失谐从负无穷扫描到正无穷、耦合常数固定且系统封闭；有限时间端点、噪声和耗散会产生Stückelberg干涉或改动转移率。
